Pour étudier la réaction entre le cuivre et le dibrome en solution, on place
dans un bécher~:
\begin{itemize}
    \item le volume $V_{1}=\qty{50}{\mL}$ d'une solution de bromure de
          cuivre(II) $\ce{Cu^2+_{(aq)} + 2Br^-_{(aq)}}$ de concentration molaire
          ${C_{1}=\qty{4e-2}{\mol\per\L}}$~;
    \item le volume $V_{2}=\qty{50}{\mL}$ d'une solution de dibrome
          $\ce{Br2_{(aq)}}$ de concentration molaire
          ${C_{2}=\qty{4e-2}{\mol\per\L}}$~;
    \item la poudre de cuivre.
\end{itemize}

\begin{donnees}
    \begin{itemize}
        \item Constante d'équilibre associée à l'équation
              $\ce{Br2_{(aq)} + Cu_{(s)} <=>[(1)][(2)]
              2Br^-_{(aq)} + Cu^2+_{(aq)}}$~: $\mathrm{K}=1,2.10^{25}$
        \item Le cuivre est en excès.
    \end{itemize}
\end{donnees}
\begin{enumerate}
    \item Identifier les couples (ox/red) intervenants dans la réaction.
    \item Calculer la valeur du quotient de réaction $Q_{r,i}$ à l'état initial
          du système chimique.
    \item Préciser, en justifiant, le sens d'évolution du système chimique.
\end{enumerate}



